\chapter{Validation en simulation}
\label{chap:validation}

\section*{Introduction}
Ce chapitre présente la validation de la chaîne complète dans Gazebo. Il décrit
l'environnement et la méthode d'essai, puis les vérifications préalables. Il expose
ensuite le décalage de repères mis en évidence par les premiers essais et sa correction,
avant de présenter les résultats mesurés après correction et les points qui restent à
valider. Toutes les valeurs données sont issues de mesures réalisées le 29~septembre~2026.

\section{Environnement et méthode d'essai}
Les essais ont été réalisés sur un poste Ubuntu~24.04 avec ROS~2 Jazzy et Gazebo Harmonic,
dans le monde \code{warehouse\_harmonic}, avec la carte \code{warehouse\_harmonic.yaml}.
La chaîne complète était lancée par \code{scripts/run-all.sh}, fenêtre Gazebo ouverte
(figure~\ref{fig:gazebo}).

\begin{figure}[htbp]
\centering
\capture{gazebo_warehouse.png}{Fenêtre Gazebo : robot AMR-X dans l'entrepôt warehouse\_harmonic.}
\caption{Robot AMR-X simulé dans l'entrepôt Gazebo.}
\label{fig:gazebo}
\end{figure}

Pour obtenir des mesures reproductibles, les scénarios de navigation ont été pilotés par
un script de test qui appelle \emph{les mêmes services} que le dashboard,
\code{/dashboard/navigation/submit} et \code{/dashboard/navigation/cancel}, et enregistre
les changements de phase publiés par la passerelle. Les destinations libres et occupées
sont tirées de la costmap globale. La vérité terrain est la pose du modèle
\code{amr\_x} lue directement dans Gazebo (\code{gz model -m amr\_x -p}). L'erreur de
position affichée est définie par
\begin{equation}
  e = \sqrt{(x_{\text{passerelle}} - x_{\text{Gazebo}})^2 + (y_{\text{passerelle}} - y_{\text{Gazebo}})^2},
\end{equation}
où la pose de la passerelle est celle transmise au dashboard, dans le repère du monde.

\section{Vérifications préalables}
Avant les essais en simulation, la compilation et les tests automatisés ont été exécutés
sous Ubuntu (tableau~\ref{tab:prealables}).

\begin{table}[htbp]
\centering\small
\begin{tabular}{@{}ll@{}}
\toprule
\textbf{Vérification} & \textbf{Résultat} \\
\midrule
Compilation \code{colcon} des paquets \code{navigation} et \code{navigation\_msgs} & réussie \\
Génération de l'interface \code{SubmitNavigationGoal} & réussie \\
Tests unitaires Python des conversions & 5/5 réussis \\
Analyse statique du frontend (\code{oxlint}) & aucune erreur \\
Tests unitaires JavaScript & 5/5 réussis \\
Compilation de production du frontend (Vite) & réussie \\
Cohérence du manifeste des mondes SDF & à jour (6 mondes) \\
\bottomrule
\end{tabular}
\caption{Vérifications préalables sous Ubuntu.}
\label{tab:prealables}
\end{table}

\section{Mise en évidence d'un décalage de repères}

\subsection{Constat}
Lors du premier essai depuis le dashboard, une destination a été choisie en
$(5{,}268 ; 2{,}448)$ sur la \emph{Live map}. Nav2 a renvoyé un succès et le dashboard
affichait le robot à destination, en $(5{,}278 ; 2{,}423)$. Pourtant, dans Gazebo, le
robot se trouvait en $(3{,}953 ; -1{,}450)$, soit à environ 4,1~m de la destination
demandée.

Pour caractériser ce défaut, la position transmise au dashboard a été comparée à la
position Gazebo en trois endroits, les décalages de repère étant nuls
(tableau~\ref{tab:avant}). L'écart est pratiquement constant, d'environ 4,1~m, et
l'erreur de cap reste inférieure à 0,6°. Il s'agit donc d'une translation pure entre les
deux repères, et non d'une erreur de localisation aléatoire.

\begin{table}[htbp]
\centering\small
\begin{tabular}{@{}lccc@{}}
\toprule
\textbf{Position} & \textbf{Pose Gazebo (m)} & \textbf{Écart $e$ (m)} & \textbf{Écart de cap (°)} \\
\midrule
Après le premier essai web & $(3{,}953 ; -1{,}450)$ & 4,094 & $-0{,}53$ \\
Après une navigation courte & $(0{,}980 ; -2{,}931)$ & 4,108 & $0{,}40$ \\
Après une annulation & $(-4{,}877 ; 1{,}458)$ & 4,131 & $-0{,}11$ \\
\bottomrule
\end{tabular}
\caption{Écart entre la position affichée et la position Gazebo, sans correction de repère.}
\label{tab:avant}
\end{table}

\subsection{Analyse}
L'emprise de la carte \code{warehouse\_harmonic.yaml} a été comparée aux limites du
monde SDF. Le monde s'étend en $x$ de $-7{,}00$ à $6{,}31$~m et en $y$ de $-9{,}93$ à
$10{,}24$~m, alors que la zone connue de la carte s'étend en $x$ de $-5{,}73$ à $8{,}27$~m et
en $y$ de $-6{,}55$ à $14{,}30$~m. La carte est donc décalée d'environ $(+1{,}3 ; +3{,}9)$~m
par rapport au monde, ce qui correspond à l'opposé de la position d'apparition du robot,
$(-1{,}2639 ; -3{,}9049)$.

Cette observation s'explique par la façon dont la carte a été produite : une carte SLAM
a pour origine la position du robot au début de l'enregistrement. L'origine du repère
\code{map} se trouve donc à la position d'apparition du robot dans le monde Gazebo
(figure~\ref{fig:reperes}). Le décalage étant nul, le dashboard dessinait des coordonnées
\code{map} sur un plan exprimé dans le repère du monde, et les destinations cliquées
étaient envoyées à Nav2 avec le même décalage.

\begin{figure}[htbp]
\centering
\begin{tikzpicture}[scale=0.48]
  \draw[gray!60,thin] (-7,-9.93) rectangle (6.31,10.24);
  \node[font=\scriptsize,gray,anchor=south west] at (-7,10.3) {limites du monde SDF};
  \draw[blue!60,thin,dashed] (-6.99,-10.45) rectangle (7.01,10.4);
  \node[font=\scriptsize,blue!70!black,anchor=north west] at (-6.99,-10.5) {emprise de la carte Nav2};
  \draw[fl,black] (0,0) -- (2.2,0) node[right,font=\scriptsize]{$x_w$};
  \draw[fl,black] (0,0) -- (0,2.2) node[above,font=\scriptsize]{$y_w$};
  \node[font=\scriptsize,anchor=south west] at (0.1,0.1) {$O_{\text{monde}}$};
  \draw[fl,blue!70!black] (-1.26,-3.9) -- (0.94,-3.9) node[right,font=\scriptsize]{$x_m$};
  \draw[fl,blue!70!black] (-1.26,-3.9) -- (-1.26,-1.7) node[above,font=\scriptsize]{$y_m$};
  \draw[red!70!black,thick,-{Stealth}] (0,0) -- (-1.2,-3.75);
  \node[font=\scriptsize,red!70!black,anchor=west] at (-0.35,-2.2) {$(t_x, t_y)$};
  \fill[red!70!black] (-1.26,-3.9) circle (4pt);
  \node[font=\scriptsize,align=center,anchor=north] at (-1.26,-4.3)
    {$O_{\text{map}}$ = apparition du robot\\$(-1{,}2639 ; -3{,}9049)$};
\end{tikzpicture}
\caption{Origine du repère \code{map} dans le monde Gazebo.}
\label{fig:reperes}
\end{figure}

\subsection{Correction}
Deux corrections ont été apportées au lancement :
\begin{itemize}
  \item la passerelle reçoit la pose mesurée de l'origine de \code{map} dans le monde :
    $t_x = -1{,}2639$~m, $t_y = -3{,}9049$~m, $\theta = 0$ ;
  \item la pose initiale envoyée à AMCL est $(0 ; 0 ; 0)$ dans \code{map}, puisque le robot
    apparaît à l'origine de la carte.
\end{itemize}

\section{Résultats après correction}

\subsection{Cohérence des positions}
Après correction, l'écart entre la position transmise au dashboard et la position Gazebo
a été mesuré en trois endroits (tableau~\ref{tab:apres}). Il reste compris entre 0,6 et
5,6~cm, avec un écart de cap inférieur à 1°. La dernière mesure a été faite environ une
seconde après une annulation en cours de déplacement.

\begin{table}[htbp]
\centering\small
\begin{tabular}{@{}lccc@{}}
\toprule
\textbf{Position} & \textbf{Pose Gazebo (m)} & \textbf{Écart $e$ (m)} & \textbf{Écart de cap (°)} \\
\midrule
Position d'apparition & $(-1{,}264 ; -3{,}905)$ & 0,014 & $0{,}33$ \\
Après la navigation B & $(-3{,}246 ; -1{,}698)$ & 0,006 & $0{,}40$ \\
Après l'annulation F & $(-0{,}016 ; -6{,}077)$ & 0,056 & $-0{,}99$ \\
\bottomrule
\end{tabular}
\caption{Écart entre la position affichée et la position Gazebo, après correction.}
\label{tab:apres}
\end{table}

\subsection{Résultats par scénario}
Le tableau~\ref{tab:resultats} présente les résultats des scénarios exécutés. Les lettres
reprennent la numérotation du protocole de validation du projet.

\begin{table}[htbp]
\centering\small
\begin{tabularx}{\linewidth}{@{}c>{\raggedright\arraybackslash}p{3.6cm}X@{}}
\toprule
\textbf{Réf.} & \textbf{Scénario} & \textbf{Résultat mesuré} \\
\midrule
A & Cohérence des positions & Écart de 0,6 à 5,6~cm et de moins de 1° en trois endroits (tableau~\ref{tab:apres}). \\
\addlinespace
B & Destination libre proche & But $(-3{,}30 ; -1{,}71)$, cap 90° : acquitté, phase \code{navigating} après 0,01~s, \code{succeeded} après 5,75~s. Erreur finale de position 5,6~cm ; cap final 101,0°. \\
\addlinespace
D & Destinations invalides & Cellule occupée et point hors carte : refusés (\og outside known free space or occupied \fg{}). Coordonnée non finie : refusée. Déplacement du robot pendant ces essais : 0,0~cm. \\
\addlinespace
F & Annulation en route & Robot en mouvement (4,1~m parcourus 6~s après l'envoi) ; annulation à 7,4~s : \code{canceling} puis \code{canceled} en 0,03~s. Déplacement dans les 3~s suivantes : 3,6~cm. \\
\addlinespace
G & Doublons et buts concurrents & Même \code{request\_id} renvoyé : refusé (\og already received \fg{}). Second but pendant la navigation : refusé (\og a goal is already active \fg{}). \\
\addlinespace
I & Simulation en pause & État publié à 3,6~Hz avant la pause, 0 message en 5~s pendant la pause, 3,8~Hz après la reprise, passerelle de nouveau prête. \\
\addlinespace
J & Mauvais monde & Demande avec \code{hospital-harmonic} : refusée (\og selected world does not match \fg{}). \\
\addlinespace
-- & Annulation sans but actif & Refusée (\og no active destination \fg{}). \\
\bottomrule
\end{tabularx}
\caption{Résultats mesurés par scénario, après correction des repères.}
\label{tab:resultats}
\end{table}

\subsection{Illustration d'une navigation réussie}
La figure~\ref{fig:succes} montre une navigation capturée simultanément dans le dashboard
et dans Gazebo. Le robot, initialement en $(-0{,}05 ; -6{,}04)$, a reçu la destination
$(-0{,}55 ; -3{,}71)$ par le service de soumission. Nav2 a renvoyé un succès après
10,8~s ; dans Gazebo, le robot s'est arrêté à 0,9~cm de la destination, et la position
affichée dans le dashboard différait de 6,2~cm de la position Gazebo. Le dashboard montre
le déplacement du robot et la trace de son trajet, et le panneau de flotte indique le
résultat \code{succeeded} (figure~\ref{fig:statut}).

\begin{figure}[htbp]
\centering
\begin{subfigure}[b]{0.34\linewidth}\centering
  \includegraphics[width=\linewidth]{nav_web_before.png}
  \caption{Dashboard, avant.}
\end{subfigure}\hfill
\begin{subfigure}[b]{0.62\linewidth}\centering
  \includegraphics[width=\linewidth]{nav_gz_before.png}
  \caption{Gazebo, avant.}
\end{subfigure}\\[0.8em]
\begin{subfigure}[b]{0.34\linewidth}\centering
  \includegraphics[width=\linewidth]{nav_web_after.png}
  \caption{Dashboard, après.}
\end{subfigure}\hfill
\begin{subfigure}[b]{0.62\linewidth}\centering
  \includegraphics[width=\linewidth]{nav_gz_after.png}
  \caption{Gazebo, après.}
\end{subfigure}
\caption{Navigation réussie : déplacement du robot vu dans le dashboard et dans Gazebo.}
\label{fig:succes}
\end{figure}

\begin{figure}[htbp]
\centering
\includegraphics[width=0.7\linewidth]{nav_status_succeeded.png}
\caption{Résultat de la navigation affiché dans le panneau de flotte du dashboard.}
\label{fig:statut}
\end{figure}

\subsection{Analyse}
Les résultats montrent que la chaîne remplit les besoins fonctionnels : la destination
est acquittée immédiatement, la phase suit fidèlement l'exécution de Nav2 et la position
affichée correspond à celle du robot dans Gazebo à quelques centimètres près.

L'écart de cap final de 11° au scénario B reste inférieur à la tolérance angulaire
configurée dans Nav2 (\code{yaw\_goal\_tolerance} de 0,25~rad, soit 14,3°) ; de même,
l'erreur de position de 5,6~cm est inférieure à la tolérance de 0,25~m. Nav2 considère le
but atteint dès que ces tolérances sont respectées ; une précision plus fine relèverait du
réglage de Nav2 et non du dashboard.

La fréquence de l'état observée en temps réel (3,6 à 3,8~Hz) est légèrement inférieure aux
4~Hz nominaux. Le temporisateur de la passerelle étant cadencé sur le temps simulé, cette
fréquence suit la vitesse d'exécution de la simulation ; c'est aussi ce qui provoque l'arrêt
de la publication lorsque Gazebo est en pause, et donc le blocage des commandes dans le
dashboard après 2,5~s.

La latence d'acquittement mesurée (2~ms) a été obtenue par un client ROS local ; elle
n'inclut pas le passage par rosbridge et le navigateur.

\section{Points restant à valider}
Certains scénarios du protocole n'ont pas encore été exécutés dans la simulation :
\begin{itemize}
  \item \textbf{C} -- détour autour d'un obstacle : non mesuré spécifiquement ;
  \item \textbf{E} -- destination libre mais inaccessible : à réaliser pour vérifier que
    l'échec de Nav2 est bien affiché ;
  \item \textbf{H} -- coupure de rosbridge, puis reconnexion sans renvoi automatique ;
  \item \textbf{K} et \textbf{L} -- mode démonstration et changement d'adresse dans
    \emph{Settings}, vérifiés jusqu'ici par le test navigateur avec un faux serveur ROS
    uniquement ;
  \item cellule \emph{inconnue} : la costmap globale utilisée ne contenait aucune cellule
    inconnue ; ce cas n'est couvert que par les tests unitaires.
\end{itemize}
Ces essais doivent être conduits depuis l'interface web, en notant pour chacun la
configuration, le but, le résultat, la durée et les observations.

\section{Limites}
La solution est conçue pour la simulation. Elle ne gère qu'un robot et un monde à la
fois, n'arbitre pas les commandes venant d'autres sources (RViz, téléopération) et
n'ajoute pas d'authentification au niveau de rosbridge : elle n'est pas destinée à être
exposée sur un réseau ouvert. Fermer la page ou perdre la connexion n'arrête pas la
navigation en cours. Enfin, l'identifiant du monde et les décalages de repère sont
configurés par l'opérateur au lancement ; ils ne prouvent pas à eux seuls que le bon monde
et la bonne carte sont chargés, d'où l'intérêt de la vérification de position décrite
ci-dessus.

\section*{Conclusion}
Les essais en simulation ont permis de valider l'envoi de destinations, le suivi de
l'état, l'annulation, le refus des demandes invalides et la cohérence de la position
affichée. Ils ont aussi mis en évidence un décalage de 4,1~m entre la carte Nav2 et le
monde Gazebo, invisible sans comparaison avec la vérité terrain, qui a été analysé et
corrigé.
